% ==============================================================================
% Section: Literature Review
% Status: EXPANDED (2026-05-02 v2) — absorbs SFE coverage that left with the
%   Model section cut; adds Anderson-Xu price-cap SFE, Hogan energy-only,
%   Hobbs RPM dynamic analysis, Vasin uniform-price SFE, Wolfram empirical
%   duopoly, Wolak/Borenstein market-power measurement, Cramton-Ockenfels
%   capacity design.
% ==============================================================================

\section{Literature Review}
\label{sec:literature}

The supply function equilibrium concept, formalized by
\citet{Klemperer1989_sfe} for demand-uncertainty oligopolies in which
each firm commits to a price-quantity schedule before demand
realizes, characterizes symmetric Nash equilibria by a first-order
differential equation whose solutions form a continuum bounded above
by the Cournot schedule and below by the competitive schedule
$p = c$. Under demand uncertainty with bounded support, every
trajectory between the two stationary solutions is itself a Nash
equilibrium, so prediction and comparative statics require an
equilibrium-selection criterion. \citet{Green1992_british} apply this
framework to the England and Wales electricity market and establish
the methodological template now standard in the calibrated SFE
literature: derive the symmetric ODE, integrate backward from the
capacity-constrained boundary, and calibrate to observed market
structure. \citet{Holmberg2008_unique_sfe} shows that in
capacity-constrained electricity markets the symmetric SFE is unique
when rivals' capacity constraints simultaneously bind with positive
probability, providing the boundary condition I use as an
equilibrium-selection device in the deterministic BRA.
\citet{Anderson2008_sfe_asymmetric} and \citet{Holmberg2007_asymmetric}
extend the existence and characterization results to asymmetric
firms; \citet{Rudkevich1998_poolco} bridges the canonical
Klemperer-Meyer theory and applied generation-market modeling.

Empirical and applied work has refined the SFE for electricity
markets. \citet{Wolfram1999_duopoly} provides the canonical empirical
test in the England and Wales duopoly and finds prices consistent
with a strategic SFE rather than the competitive limit;
\citet{Sweeting2007_market_power} extends the analysis through
restructuring. \citet{Sioshansi2007_sfe_ercot} fits the SFE to ERCOT
balancing-market data and assesses its predictive performance.
\citet{Baldick2004_linear_sfe} develop a tractable
linear-affine variant useful for transmission-constrained networks,
and \citet{Vasin2016_sfe_uniform} characterize the SFE specifically
for uniform-price auction mechanisms of the kind PJM operates.
\citet{Anderson2005_sfe_pricecaps} is the most direct theoretical
antecedent of the present paper: they derive supply function
equilibria in spot markets with contracts and price caps and show
that a binding cap pushes the equilibrium toward the capacity-cap
boundary, with quantity rationed at the cap when supply is
insufficient. The settlement studied here is one practical
realization of that mechanism.

The capacity-market design literature originates from the
``missing-money'' problem: regulatory price caps in the wholesale
energy market suppress the scarcity rents that would otherwise
finance investment in peaking generation, leaving energy revenues
insufficient to recover fixed costs of new entry
\citep{Joskow2007_capacity, Joskow2008_capacity_payments}.
\citet{Cramton2005_capacity} and \citet{CramtonOckenfels2012_capacity}
develop the design principles for a forward capacity instrument with
a sloped administrative demand curve, and
\citet{Cramton2007_uniform} defends uniform-price auctions over
discriminatory alternatives in this setting. PJM's Reliability
Pricing Model is analyzed dynamically by \citet{Hobbs2007_rpm}, whose
proposal informed the design of the VRR curve, and
\citet{Bowring2013_pjm} provides the institutional account from the
Independent Market Monitor's perspective; \citet{Pfeifenberger2011_rpm}
offers an early performance assessment. \citet{Hogan2005_energy_adequacy}
articulates the energy-only alternative, in which scarcity prices in
the wholesale energy market are allowed to clear at administrative
ceilings calibrated to the value of lost load and so deliver the
investment signal that capacity-market advocates argue is otherwise
missing.

The economics of binding price ceilings is foundational.
\citet{Weitzman1974_prices} shows that the choice between price and
quantity instruments depends on the relative slopes of marginal
benefit and marginal cost; when supply is inelastic in the short
run, price instruments dominate. \citet{GlaeserLuttmer2003_rent}
demonstrate, in the rent-control context, that binding ceilings
generate a distinctive consumption misallocation across
heterogeneous demanders, qualitatively different from the standard
deadweight triangle. Experimental evidence in
\citet{Vossler2009_price_caps} confirms the qualitative shortage
prediction in laboratory uniform-price wholesale-electricity
markets. The cap interventions this literature has primarily
studied are spot-market bid caps and offer-level mitigation rules
of the kind PJM's Three Pivotal Supplier test imposes; the Shapiro
settlement is distinct in being a prospective, market-wide cap on
a forward capacity instrument, which is part of why the textbook
shortage prediction has such a clean realization in this setting.
The empirical market-power literature on restructured electricity
markets---\citet{Wolak2003_measuring},
\citet{Borenstein2002_market_power}, and
\citet{Bushnell2008_vertical}---supplies the methodological backdrop
for distinguishing competitive from strategic clearing and motivates
the cap-incidence diagnostic of
Section~\ref{sec:cap_incidence}. The applied question of this paper
is the realization of these classical predictions in PJM's capacity
market under the Shapiro settlement, with attention to the
investment-signal cost of a binding cap during a tight-supply
episode \citep{Cramton2005_capacity}.
